\chapter{Central Claim-Status and Open-Obligation Ledger}
\label{app:claim-ledger}

This appendix collects the manuscript's unresolved obligations together with
the resolved, refuted, and conditional boundaries needed to state the surviving
scope.  It does not upgrade any claim beyond its local theorem statement or
certificate.  A \emph{conjecture} is a favored, falsifiable statement that has
not been proved; an \emph{open problem} has no presumed answer.  The local
theorem statements, certificates, and hypotheses remain authoritative.

\small

\section{Mathematical and probabilistic closure}

\begin{description}
\item[Continuum law theory (open).]
Construct a refining projective family of normalized finite laws, prove
tightness in a declared topology on sections, control gauges and reference
measures, and pass the ELBO, connections, and observables to the limit.
Finite-dimensional consistency alone does not do this. \status{OPEN}

\item[Pointwise-dominated and nondominated likelihood families (open).]
The finite interaction theory uses one family-level dominating measure for
each record kernel and a finite jointly measurable density version.  If
domination exists only separately at each latent state, pointwise
Radon--Nikodym versions need not supply a common measurable likelihood
$L_o(y)$.  Extend the construction directly at the joint-measure and regular
conditional-law level, or prove hypotheses giving jointly measurable versions,
and specify how record-null version choices propagate through local
disintegrations and coarse maps. \status{OPEN}

\item[Regular frame-coordinate quotient (open).]
Determine generic stabilizers of the passive \(G\times G\) action generated by
independent belief and model trivializations, and identify strata on which the
orbit space is a regular quotient. This coordinate action is not the
principal structure group. A dimension subtraction is invalid without
freeness and properness. \status{OPEN}

\item[Optimization, scheduling, and projection (open).]
The finite conditional-gap and local--global potential identities establish
objective compatibility, not convergence.  For each proposed algorithm,
declare the sequential, asynchronous, or parallel schedule and give sufficient
compact-sublevel, lower-semicontinuity, differentiability, and exactness or
line-search/damping hypotheses for convergence to a stationary point, or give
a counterexample in the declared family.  Parallel replacement of correlated
full conditionals additionally requires proof that the replacements define one
joint law.  Separately, prove when Bregman projection is gauge compatible and
commutes with a coarse parameter map, and determine whether a positive mass
pin is relevant, marginal, or irrelevant under that map. \status{OPEN}

\item[Joint-law lift and selector boundary (open/refuted split).]
An exact VFE history requires the commuting conditional/full-law lift of
\Cref{hyp:hist-exact-vfe-lift}, including the dependence coordinates omitted
by marginal belief--model sections.  Existence of such a smooth right inverse,
nondegeneracy of its Fisher pullback, and preservation under coarse graining
must be proved for each declared recognition family; they do not follow from
the two associated bundles alone.  The governing finite correction in
\Cref{thm:hist-three-agent-parity-lift} closes one smooth positive categorical
instance: one parity direction inside the 57-dimensional Frechet fiber for
three paired binary agents.

The governing selector correction of 2026-08-14 supersedes the unqualified
claim that canonical selection is merely open.  On finite typed lists,
naturality under coordinatewise finite Markov kernels whose randomness is
independently tensored uniquely forces the product section; this quantifier
does not cover all of \(\mathsf{FinStoch}\).  A wide marginal-compatible
enlargement containing that local class admits a natural section exactly when
every added morphism preserves products.  Separately, the frozen contract
explicitly admits both split kernels \(R_{1/3}\) and \(R_{1/2}\).  From one
fair source they have the same fair/fair target marginals but respective atom
multisets
\(\{1/3,1/3,1/6,1/6\}\) and
\(\{3/8,3/8,1/8,1/8\}\).  Naturality under both declared splits would force
one single-valued target section to return both unequal laws, even modulo
relabeling.  The admission of both splits and single-valuedness are
load-bearing.  The absolute selector with those quantifiers is
\emph{REFUTED}.

Under the declared full naturality and canonicity requirements, selecting
nonproduct dependence requires either restricting the admitted morphism class
or adding declared relational, constraint, or reference data.  Relative to a
finite reference \(p\), statistic \(T\), and
target \(m\), a finite-KL minimizer of \(\KL(q\Vert p)\) subject to
\(\E_qT=m\) exists exactly when
\(m\in\operatorname{conv}T(\operatorname{supp}p)\) and is unique as a law.
For deterministic \(f:X\to Y\), conditional-KL completion is defined for
\(r\ll f_\#p\) and composes along nested coarse maps when each stage uses
the pushed reference.  These relative constructions do not close
arbitrary-family existence, coarse preservation, odd independent relabelings,
\(\GL(K)\) equivariance, a continuum lift, or renormalization.
\status{OPEN}

\item[Operational intervention quotient and nonidentifiability (established).]
Fix a monoid \(A\) and a response \(\Phi\).  Contextual equivalence is the
largest congruence contained in the kernel relation of \(\Phi\).  For a
compatible quotient \(q\), write
\[
\begin{gathered}
\pi:A\twoheadrightarrow\operatorname{Syn}(\Phi),\qquad
q:A\twoheadrightarrow B,\\
\Phi=\psi q.
\end{gathered}
\]
There is one unique surjective unital homomorphism satisfying
\[
\begin{gathered}
h:B\twoheadrightarrow\operatorname{Syn}(\Phi),\\
\pi=hq,\qquad \bar\Phi h=\psi.
\end{gathered}
\]
For finite \(A\), this gives minimum
protocol-class cardinality and uniqueness over \(A\), not a minimum raw
DAG or latent realization.  Under compact-metrizable monoid, continuous
response, and metrizable-Hausdorff target hypotheses, a countable dense
context signature realizes a compact metrizable quotient with continuous
multiplication and response; see
\Cref{prop:hist-operational-quotient-universal-property,thm:hist-compact-operational-quotient}.

The declared groupoid
\(\mathsf{FinTIP}_{R,O}^{\mathrm{iso}}\) contains finite normalized typed
DAG presentations and all partial hard assignments.  Total right override
forms the intervention monoid, and equality of every two-sided contextual
retained response is a congruence.  Thus
\(\operatorname{Red}:\mathsf{FinTIP}_{R,O}^{\mathrm{iso}}
\to\mathsf{FinRIE}_{R,O}^{\mathrm{iso}}\) and
\(\overline U_{\mathrm{pass}}:\mathsf{FinRIE}_{R,O}^{\mathrm{iso}}
\to\mathsf{FinObs}_{R,O}^{\mathrm{iso}}\) are well-defined.

Within the declared BSC subcategory,
\(L(1/4,1/3)\) and \(L(1/3,1/4)\) have the same passive retained law but
nonisomorphic fifteen-class hard experiments.  The independent null extension
collapses and is a control.  Consequently no universal
\(R\overline U_{\mathrm{pass}}\cong\operatorname{id}\) recovers every
reduced object, while
\(\overline U_{\mathrm{pass}}R\cong\operatorname{id}\) may still select
one representative.

The same pair has unequal exact marked-soft mediator-face diameters,
\((1-2\epsilon)/3\) and \((1-2\epsilon)/2\), for
\(0<\epsilon<1/2\), including strict-interior witnesses.  For independent
affine randomization, a fifteen-coordinate contextual minor has determinant
\((2b-1)^6(2\delta-1)^3/32\), so any admitted randomized isomorphism would
restrict on simplex vertices to the forbidden hard isomorphism; see
\Cref{thm:hist-soft-bsc-target-face-nonidentifiability,thm:hist-randomized-hard-intervention-nonidentifiability}.
The randomized proof does not reuse the hard unmatched-response invariant,
which convexification destroys.

For a finite DAG of standard-Borel node spaces and palettes, declared
normalized pointwise kernel families with jointly measurable evaluations give
a Borel retained response.  The construction supplies an algebraic contextual
quotient but does not by itself establish a standard-Borel quotient; that
requires an exhibited smooth classifier or stronger topology.  Compact-Polish
node spaces, compact palettes, isolated baseline symbols, and jointly Feller
kernels give a compact metrizable quotient with weakly continuous response.
The circle heat pair has equal passive retained law but a strict one-way Blackwell
order and strict soft response-set inclusion under the retained mediator,
heat geometry, ordered boundary, and one global response map; see
\Cref{prop:hist-standard-borel-intervention-semantics,cor:hist-compact-feller-operational-quotient,thm:hist-circle-heat-intervention-nonidentifiability}.
\status{ESTABLISHED}

\item[Category-independent intervention semantics (open).]
The released extension supersedes blanket open wording only for its fixed
hard, normalized marked-soft, independently randomized, declared
standard-Borel, compact-Feller, and circle categories.  Target-erasing soft
maps, correlated or adaptive randomization, and identification of null-version
point interventions from almost-sure passive observational conditionals remain
open, as do noncompact quotients, category canonicity, arbitrary latent
dilation, and minimal raw realization.  No result establishes fixed-observation
ELBO or VFE equality, autonomous agency, continuum base-manifold/gauge/RG
dynamics, or any identification with physical geometry, time, units,
constants, or ontology; see
\Cref{open:hist-typed-intervention-recovery}. \status{OPEN}

\item[Global relational information clock (open).]
The finite-dimensional Fisher length of one regular oriented history is
defined and reparameterization invariant.  A scalar clock shared across
agents or patches additionally requires the exactness and period conditions
isolated in \Cref{thm:hist-global-clock-exactness}, compatibility across
critical and null strata, and, for a continuum of sections, the missing base
measure and function-space theory.  None of those follow from local
dissipation alone. \status{OPEN}
\end{description}

\section{General coarse maps and renormalization}

\begin{description}
\item[Full pointwise probabilistic datum (established).]
Fix a nonempty finite child block \(I\), a parent label \(A\), one
\(r_*\in\mathcal U_A=\bigcap_{i\in I}\mathcal C_i\), and structural data \(X\)
with \(X_A=\chi_A(X)\) outside the random channel. Under the hypotheses of
\Cref{thm:rg-pointwise-parent-datum}, one normalized recognition-independent
\(C_A:\mathsf Y_I\rightsquigarrow\mathsf Z_A\) pushes the fixed fine generative
joint, selected posterior-version family, and correlated recognition law to a
normalized parent triple while leaving the observation in \(\mathsf O\)
unchanged. An induced parent evaluator exists by standard-Borel disintegration;
a predeclared evaluator requires separate almost-sure compatibility. All named
recognition, prior, and posterior marginals are projections of their typed full
laws and do not reconstruct the corresponding joints.

The KL chain has the nonnegative conditional-information defect
\(\Delta_A(o,X)\) as an additive identity in \([0,+\infty]\). Adding the same
finite real \(-\log p_X(o)\) to both KL terms gives an extended-real VFE
identity; a finite VFE may be negative. Without a finiteness premise,
\(\Delta_A=0\) exactly when the discarded conditional recognition and
posterior laws agree \(\mathbb Q_{A,o,X}\)-almost surely. Finite fine KL is
required for ordinary subtraction
\(\Fenergy_I-\Fenergy_A=\Delta_A\) and for the stated two-way pairwise
common-recovery equivalence. Holonomy blindness requires
typed actions, full fine-law covariance, compatible selected posterior versions,
\(C_A\) equivariance, and evaluator covariance; same-slice invariance requires
fixed-\((o,X)\) isotropy. Raw holonomy retention is the alternative and makes no
blindness or membership-selection claim. The manuscript status is
\textsc{established}; the bound release records ledger state
\texttt{EVIDENCE\_VERIFIED} and terminal package status
\texttt{COMPLETE\_AFFIRMATIVE}. \status{ESTABLISHED}

\item[Partition selection and experiment-level recovery (open).]
\Cref{thm:cg-holonomy-kl-marginal} proves that the marginal-law distortion is
invariant under gauge, path, and root choices. It characterizes its attained
zero as a holonomy-stabilized parallel marginal section.  The normalized
evidence-preserving full pointwise joint-law pushforward and finite pairwise
recovery criterion are supplied by \Cref{thm:rg-pointwise-parent-datum}; what
remains open is one common recovery kernel for an entire statistical experiment
without separately imposing simultaneous family hypotheses, and a relabeling-natural,
gauge-compatible nondegenerate partition selector. \status{OPEN}

\item[Downstream comparison, gluing, and agency (open).]
The next order is fixed but none of its members lies in the static release
ancestry. First, a comparison theorem must freeze the parent experiment and a
typed category for target retention, the ordered \(R\to E\to O\) boundary,
time orientation, and protocol-independent relabeling; interventions remain
analyst probes. Second, a patchwise theorem over \(\mathcal U_A\) must supply
measurable or smooth channel families, cocycles, changing active components,
strata, normalized soft multiple membership, holonomy selection, and an
integrable defect. Only that theorem can construct parent local sections and a
geometric meta-agent. Third, participatory nonequilibrium and emergent agency
require one tower action or proved reduction without double counting, a typed
choice among gradient, open, antisymmetric, kinetic, stochastic, delayed,
memory, or adaptive mechanisms, and operational agency grades. Canonical
channel and membership selection, autonomy, physical time, continuum limits,
unique microscopic physics, and ontology remain open. \status{OPEN}

\item[Infinite-volume RG limit (open).]
The finite exact scale theory supplies normalized kernels, pushed measure
pairs, a local analytic bounded-action calculus, its \(L^p\) conditional
expectation extension, the exact conditional-variance loss, full
power-set Hoeffding interaction coordinates, and an exact nonlinear
interaction map with a typed projection residual.  The interaction theorem
separately assumes an equivalent product reference at every admitted scale;
an arbitrary normalized Markov arrow need not preserve that tier.  None of
these finite-step theorems supplies a thermodynamic limit.
An infinite-volume limit still requires projectively consistent cylinder laws,
boundary conditions or a DLR specification, tightness in the selected state
topology, convergence of free-energy densities, and compatible limits of the
connection and cross-morphism components. \status{OPEN}

\item[Two-index limits and universality (open).]
Prove a nontrivial diagonal limit or an exchange theorem for system size
\(n\) and RG depth \(\ell\).  Then establish convergence of invariant
observables, basin independence, and blocking-scheme robustness.  A finite
congruence identity or one matched exponent is insufficient. \status{OPEN}

\item[Extensive relevance beyond the scalar Gaussian realization (open).]
\Cref{thm:rg-gaussian-hermite-spectrum} inhabits all three relevance classes for
independent scalar Gaussian coordinates under the declared replication sum,
block statistic, and fixed Fisher norm.  \Cref{prop:rg-hermite-scope} shows that
correlated coordinates, unrestricted multivariate covariance, and general
\(\GL(d)\) actions each leave that theorem.  Obtaining the corresponding
spectrum for a correlated, anisotropic, or gauge-mixed block requires a declared
transported family of laws and tangent norms together with a separate
diagonalization. \status{OPEN}

\item[Nonlinear attraction at the Gaussian tangent (open).]
The Hermite spectrum controls the derivative of the score map at the Gaussian
law.  A basin of attraction, a remainder estimate for the nonlinear map, and the
fate of the marginal degree-two direction do not follow from a linear spectral
statement and are not supplied here. \status{OPEN}

\item[Multiplicative ergodic structure of the interaction cocycle (open).]
\Cref{def:rg-mode-line} and \Cref{thm:rg-tempered-comparison} define compatible
cross-scale modes and give exponent invariance under tempered comparison data.
An Oseledets splitting would additionally require a measure-preserving base
flow, measurability of the cocycle in the base variable, integrability of
\(\log^+\lVert D\rVert\) and of \(\log^+\lVert D^{-1}\rVert\), and either finite
dimension or a compactness hypothesis.  A deterministic blocking sequence
supplies none of these, so the growth rates remain definitions rather than an
existence theorem. \status{OPEN}

\item[Smooth scale tier and continuous beta (open).]
\Cref{def:rg-scale-connection} declares the scale manifold, Banach bundle, and
scale connection that a continuous beta requires, and
\Cref{prop:rg-continuous-beta-underdetermined} shows that a discrete blocking
sequence does not determine them.  Selecting a canonical interpolation, or
proving that a reported continuous beta is independent of the choice inside a
declared class, remains open. \status{OPEN}

\item[Nontrivial invariant retained subspace (open).]
\eqref{eq:rg-beta-exactness-criterion} settles exactly when a retained beta
function is exact.  Whether any retained subspace other than the zero and full
ones is invariant under \(T_\ell^{\mathcal G}\) for an admitted network and
blocking channel is not settled here, so every reported projected flow carries
its residual \(\delta\beta_\ell\). \status{OPEN}

\item[Bayesian-RG bridge (open).]
To compare this flow with Bayesian renormalization, place both on one state
space, transport the relevant Fisher tensors, declare a common scale map,
and prove a monotone.  Similar Gaussian formulas do not prove opposite flow
directions. \status{OPEN}

\item[Fine--coarse history semiconjugacy for the declared RG maps (open).]
\Cref{thm:hist-record-clock-contraction} compares Fisher duration along a
single path pushed through a parameter-independent record channel, and
\Cref{thm:hist-oriented-semiconjugacy,prop:hist-noncollapse,thm:hist-duration-relation,thm:hist-duration-criterion}
settle what oriented semiconjugacy does and does not give once it holds.  What
is not proved is that this manuscript's own independently recomputed coarse
natural-gradient flow satisfies it.  One structured sufficient route is the
one isolated at \Cref{prop:hist-natural-gradient-sufficiency}: exhibit
\(\chi_\ell\) with \(\chi_\ell'>0\) and an open set containing the fine orbit on
which \Cref{hyp:hist-functional-compatibility} holds, prove for the
exact-contraction coarse objective that the orbit lies in the interior of the
attaining set, and then verify horizontal conformality of the declared
\(\mathsf R_\ell\).  This route is not asserted necessary or minimal.  Behavior
at critical and singular strata remains part of the open work.  The route is
satisfiable, since the exhibited tier of
\Cref{def:hist-finite-configuration-tier} meets it with
\(a_\ell=\tfrac12\). \status{OPEN}

\item[Configuration tier for the declared recognition family (open).]
\Cref{def:hist-finite-configuration-tier} and
\Cref{thm:hist-finite-tier-regularity} exhibit a nonempty finite-dimensional
configuration manifold with a strong constant Gram metric, a rank test for
nondegeneracy, and a closed gauge-orbit tangent, and the infinite-dimensional
\(L^2\) and Sobolev tiers are separated with the failure witness attached.
What is not settled is whether the manuscript's own declared recognition family
can be brought to one of those tiers at every scale at which a natural-gradient
field or a Fisher duration is asserted.  Until that is done, the history
theorems hold on the exhibited tier and are conditional elsewhere.
\status{OPEN}

\item[Smooth structure of the projectable set (open).]
\Cref{thm:pb-section-descent} characterizes which fine sections descend, and
\Cref{prop:pb-nonfunctorial-descent} shows the descendable set is proper and can
have empty interior.  Whether that set is a smooth submanifold of a declared
configuration manifold carrying a smooth induced arrow needs a transversality
or elliptic-regularity hypothesis that is not supplied.  This manuscript avoids
the question by declaring \(\mathsf R_\ell\) separately; the pointwise-induced
route remains open for anyone who prefers it. \status{OPEN}

\item[Proper, free, and isometric gauge action in infinite dimensions (open).]
The available witness for failure of the orthogonal quotient speed is free and
isometric with a dense, hence nonclosed, orbit tangent.  Properness would force
closed orbits, so that witness does not settle whether free, proper, and
isometric action alone suffices.  Either a proper witness or a proof under
properness is required. \status{OPEN}
\end{description}

\section{Multivariate-Gaussian realization}

\begin{description}
\item[Scalarized attraction (conjecture).]
\Cref{conj:grg-fixed-b-attraction} predicts a unique attracting spatial ray
only after the matrix direction \(M_0\) and one primitive spatial map \(B\)
are fixed and the Perron vector factorizes.  Primitivity, factorization, and
the fixed hierarchy still require proof in any application.  The claim does
not cover the full matrix cone or a nonautonomous cocycle.
\status{CONJECTURE}

\item[Admissible cone classification (open).]
Extend the exact classification beyond the fixed
congruence-diagonal cone to correlated or nonconvex constraints, variable
charts, parameter-dependent domains, and nonflat link data. \status{OPEN}

\item[Nonflat-link compression (open).]
Characterize when the rectangular endpoint data of a blocked cut
compress to one invertible transporter and one SPD weight, remains closed
under another merge, and admits a normalized probability and rescaling lift.
\status{OPEN}

\item[Intrinsic scale selection (open).]
Find a nondegenerate criterion internal to one fixed evidence problem, or
prove endpoint degeneracy for a stated class.  A successful alternative may
be a normalized Markov pushforward or a proved Lyapunov functional; the
determinant gap and quotient-volume convention do not supply one. \status{OPEN}

\item[Information-geometric transfer (open).]
Identify the coarse recognition parameter with a parameter of the fitted
generative law, prove nondegeneracy of the relevant Fisher metric, and then
test the proposed scale construction.  Recognition-space and model-space
Fisher tensors are not interchangeable by name. \status{OPEN}

\item[Stochastic inverse RG (open).]
Choose compatible priors for discarded internal edges, nested refinement
kernels, chart data, and a projective-limit theorem.  Infinite divisibility
can redraw compatible microstates in distribution; it cannot invert a
many-to-one realized blocking. \status{OPEN}

\item[Update robustness (open).]
Extend the anchored Gaussian convergence theorem, or find its boundary, for
delayed, noisy, asynchronous, or inexact updates.  The present contraction
proof covers only its stated synchronous exact regime. \status{OPEN}
\end{description}

\section{Geometric and physical interpretation}

\begin{description}
\item[Operational base holonomy (open).]
Choose a named principal bundle and connection, an assigned base loop, a
gauge-invariant population-record statistic, and a typed map from connection
data to its record law; then construct two admissible connection data sets
with every declared non-holonomy input fixed, distinct represented loop-holonomy
conjugacy classes, and differing statistic laws, or prove insensitivity for
that same declared controlled class.  A stronger sufficient formulation
requires the record-law map to factor nontrivially through the represented
holonomy conjugacy class.
Graph-link holonomy alone does not settle this question, and no such
operational tuple or empirical evidence is presently supplied. \status{OPEN}

\item[Graph-to-base identification (open).]
Embed the interaction complex in \(\mathcal C\), assign oriented
piecewise-smooth curves, and prove
\(\widehat\rho_b(\Theta_e^b)=\Omega_{\gamma_e}\) and
\(\widehat\rho_m(\Theta_e^m)=\widetilde\Omega_{\gamma_e}\).
Even then the result concerns connection holonomy, not bundle topology.
\status{OPEN}

\item[Canonical nondegenerate pullback geometry (open).]
\Cref{thm:pb-pullback-gauge-invariance} proves passive gauge covariance for a
fixed section and chosen connection, while
\Cref{thm:pb-pullback-rank-quotient} identifies the conditional quotient
geometry.  A canonical or empirically specified connection, relative channel
weights, constant rank, integrability of the radical, a regular Hausdorff leaf
space, basicness of both induced tensors, and a physical interpretation remain
separate obligations.  No canonical connection is selected anywhere in this
manuscript.  Every pullback tensor, and every mixed or section-induced length
that uses a horizontal--vertical split, stays connection relative; every
Fisher duration stays relative to its declared metric.  By contrast, the
vertical Fisher length of a curve confined to one fixed statistical fiber uses
only the vertical Fisher metric and is connection independent, although it
remains metric relative. \status{OPEN}

\item[Retained Fisher quotient and declared blocks (conditional boundary).]
For a smooth retained-law map \(\rho:\Theta\to N\) and a positive-semidefinite
target tensor \(g\), the exact pullback radical is
\[
\operatorname{rad}(\rho^*g)
 =d\rho^{-1}(\operatorname{rad}g).
\]
It equals \(\ker d\rho\) exactly when
\(\operatorname{im}d\rho\cap\operatorname{rad}g=\{0\}\).  Under this
transversality and constant rank, the visible tangent is a smooth
positive-definite quotient vector bundle.  A global quotient manifold still
requires a simple Hausdorff null-leaf quotient.  Under the equality
\(\operatorname{rad}(\rho^*g)=\ker d\rho\) and a regular connected-leaf
quotient, the specific tensor \(\rho^*g\) is basic because \(\rho\) factors
through that quotient.  This does not make an arbitrary tensor, a
declared block projector, or the connection-relative covariant jet of
\Cref{hyp:pb-regular-models} automatically basic.

For a declared tangent splitting, direct quotient blocks require blockwise
kernel splitting; additive Fisher energy additionally requires pairwise
orthogonality; and block descent requires projectable/basic image data,
including the global disconnected-fiber condition when applicable.  In the
promoted six-bit parity family the full-joint map has rank seven everywhere,
whereas singleton retention has rank six with kernel exactly
\(\operatorname{span}\{\partial_\kappa\}\).  This rank statement neither
supplies node typing nor agentizes \(\kappa\).  The vector-bundle identities
are established under their stated hypotheses; a global regular quotient,
canonical block data, autonomous agency, and physical interpretation remain
open. \status{OPEN}

\item[Necessity in the base Fisher cocycle and the positivity criterion (open).]
\Cref{thm:pb-base-defect-cocycle} gives the exact residual of the base
Fisher-defect cocycle and shows that the sharp cocycle holds exactly under the
quadratic-form equality \eqref{eq:pb-base-cocycle-criterion} (an equal-seminorm
criterion only under the Markov/positive-semidefinite hypothesis), for which
three sufficient conditions are recorded and none is necessary.  Similarly,
\Cref{thm:pb-signed-positivity-criterion} shows that a vanishing horizontal
defect is sufficient and not necessary for base positivity.  Classifying the
accidental solutions of either criterion, in terms of intrinsic data rather
than by inspection of a chart, is open. \status{OPEN}

\item[Physical-time identification (open).]
Fisher duration on an oriented inference history is unchanged by an
orientation-preserving change of parameter.  It is a statistical length, not
a physical clock.  A physical-time claim would require a named target clock,
a calibration map and uncertainty model, causal and compositional consistency,
and discriminating evidence against alternative parameterizations.  None is
supplied here. \status{OPEN}
\end{description}

\openproblemheading{Physical-law identification}{claim:physical-law-identification}
Identifying a repaired limiting object with a physical law requires a named
target system, an observable and estimator, a frozen validation and
uncertainty protocol, a baseline margin, probability-level closure, and
held-out evidence robust across admissible blocking schemes.  None is
supplied here.  Accordingly, the present mathematical theory neither proves
a physical law nor supports ontological closure. \status{OPEN}

\normalsize
